\documentclass[aspectratio=169,10pt,english]{beamer}

\input{../../../_preambulo_beamer.tex}
\input{../../../_comandos_beamer.tex}

\selectlanguage{english}

\title{MA1001B - Statistical Modeling for Decision Making}
\subtitle{Lesson 02: Probability Sampling Designs}
\author{}
\institute{Tecnol\'ogico de Monterrey}
\date{}

\begin{document}

\begin{frame}[plain]
  \vspace{1.2cm}
  {\Large\bfseries MA1001B - Statistical Modeling for Decision Making\par}
  \vspace{0.7cm}
  {\large Lesson 02: Probability Sampling Designs\par}
  \vspace{0.7cm}
  {\color{TecRojo}\rule{\textwidth}{1.2pt}}
  \vspace{0.8cm}

  MA1001B Faculty\\[0.3cm]
  Term\\[0.3cm]
  Tecnol\'ogico de Monterrey
\end{frame}

\begin{frame}{The Sampling-Design Question}
  \begin{alertblock}{Guiding question}
    Which units will be selected at random, and which parts of the population
    will the resulting sample cover?
  \end{alertblock}

  \vspace{0.5em}
  By the end of this lesson, you can:
  \begin{itemize}
    \item connect a probability sample to a sampling frame;
    \item implement simple random and proportional stratified sampling;
    \item identify the selected unit in one-stage cluster sampling;
    \item compare designs without declaring a universal winner.
  \end{itemize}

  \vfill
  \scriptsize All records are synthetic. No real company or facility data are used.
\end{frame}

\begin{frame}{Three Probability Sampling Designs}
  \small
  \begin{center}
    \begin{tabular}{p{0.20\textwidth}p{0.27\textwidth}p{0.39\textwidth}}
      \toprule
      \textbf{Design} & \textbf{Randomized unit} & \textbf{Construction} \\
      \midrule
      Simple random & Individual order & Select $n$ orders from one complete frame \\
      Stratified & Order within a region & Partition into strata; sample from every stratum \\
      One-stage cluster & Facility & Select facilities; include every order inside them \\
      \bottomrule
    \end{tabular}
  \end{center}

  \vspace{0.7em}
  \begin{exampleblock}{Synthetic frame}
    $12{,}000$ orders $=$ 24 facilities $\times$ 500 orders. Four regions define
    strata; facilities define clusters.
  \end{exampleblock}

  \textbf{Probability design:} selection is governed by a known chance process,
  not by which records happen to be easiest to reach.
\end{frame}

\begin{frame}[fragile]{Step 1: Simple Random Sampling}
  \begin{columns}[T]
    \begin{column}{0.42\textwidth}
      \small
      \begin{lstlisting}
sample = population.sample(
    n=240,
    replace=False,
    random_state=42,
)

inclusion_chance = n / N
      \end{lstlisting}

      \begin{block}{Verified output}
        $N=12{,}000$, $n=240$\\
        Inclusion chance $=2.00\%$\\
        $\mu=47.67$, $\bar{x}=47.17$ min\\
        Absolute error $=0.50$ min
      \end{block}
    \end{column}
    \begin{column}{0.58\textwidth}
      \centering
      \includegraphics[width=\textwidth]{../figures/sampling_designs_01_simple_random.png}
      \vspace{0.15em}
      \scriptsize Equal selection chance does not force exact regional balance.
    \end{column}
  \end{columns}
\end{frame}

\begin{frame}[fragile]{Step 2: Proportional Stratified Sampling}
  \begin{columns}[T]
    \begin{column}{0.40\textwidth}
      \small
      \textbf{Allocation rule}
      \[
        n_h = n\frac{N_h}{N}
      \]
      Here $N_h=3{,}000$ for every region, so $n_h=60$.

      \vspace{0.4em}
      \begin{lstlisting}
for region in regions:
    stratum = population[
        population.region == region
    ]
    sample 60 orders
      \end{lstlisting}

      SRS error: 0.50 min\\
      Stratified error: 0.28 min
    \end{column}
    \begin{column}{0.60\textwidth}
      \centering
      \includegraphics[width=\textwidth]{../figures/sampling_designs_02_stratified.png}
      \vspace{0.15em}
      \scriptsize Stratification guarantees representation, not the smallest
      error in every possible sample.
    \end{column}
  \end{columns}
\end{frame}

\begin{frame}{Step 3: One-Stage Cluster Sampling}
  \begin{columns}[T]
    \begin{column}{0.40\textwidth}
      \small
      Randomly selected facilities:
      \[
        \{\mathrm{F02},\mathrm{F11},\mathrm{F16},\mathrm{F18}\}
      \]
      Include all 500 orders from each selected facility.

      \vspace{0.5em}
      \begin{tabular}{lr}
        \toprule
        Orders included & 2,000 \\
        Inclusion chance & 16.67\% \\
        Regions present & 3 of 4 \\
        Mean estimate & 48.76 min \\
        Absolute error & 1.09 min \\
        \bottomrule
      \end{tabular}

      \vspace{0.5em}
      Nearby units can share a facility effect; a large cluster sample may carry
      less independent information than its row count suggests.
    \end{column}
    \begin{column}{0.60\textwidth}
      \centering
      \includegraphics[width=\textwidth]{../figures/sampling_designs_03_cluster.png}
      \vspace{0.15em}
      \scriptsize Pink bars mark the selected facilities in this seeded sample.
    \end{column}
  \end{columns}
\end{frame}

\begin{frame}{Choosing a Design}
  \small
  \begin{center}
    \begin{tabular}{p{0.18\textwidth}p{0.23\textwidth}p{0.22\textwidth}p{0.22\textwidth}}
      \toprule
      \textbf{Design} & \textbf{Frame needed} & \textbf{Coverage} & \textbf{Operational tradeoff} \\
      \midrule
      Simple random & Complete order list & Random, not fixed by region & Units may be scattered \\
      Stratified & Order list plus region & Every region sampled & Needs region labels \\
      Cluster & Facility list & Selected facilities only & Fewer sites to reach \\
      \bottomrule
    \end{tabular}
  \end{center}

  \vspace{0.8em}
  \begin{alertblock}{Interpret the comparison correctly}
    The observed errors---0.50, 0.28, and 1.09 minutes---belong to one synthetic,
    seeded realization. They illustrate consequences of population structure;
    they do not rank the three designs for every setting.
  \end{alertblock}
\end{frame}

\begin{frame}{Concept Check}
  \begin{enumerate}
    \item In the SRS, what is one order's inclusion chance, and what frame is
      required to implement the design?
    \item Why does the stratified sample contain every region while the SRS does
      not guarantee exact regional counts?
    \item In the cluster design, what is selected at random, and which orders
      enter the sample afterward?
    \item Why does this seeded comparison not prove that stratification is
      always the most accurate design?
  \end{enumerate}

  \vspace{0.5em}
  \begin{block}{Connection to Lesson 03}
    Once a defensible sample has been collected, numerical summaries describe
    its center and variability.
  \end{block}

  \vfill
  \scriptsize Source: OpenStax, \textit{Introductory Business Statistics 2e},
  Chapter 1, Section 1.2. CC BY-NC-SA 4.0.
\end{frame}

\appendix



\end{document}
